\begin{remark}\label{remark-radical-definition-and-proof}
The \wikilink{a2s29--b001}{radical definition} and \wikilink{a2s29--b002}{intersection-of-primes theorem} apply to arbitrary commutative unital rings, not just polynomial rings over algebraically closed fields. Taking the zero ideal gives the nilradical statement; this is already an instance of the existing theorem, not an additional theorem requiring a new proof.

Read the source proof together with this section's correction note: the printed argument does not justify the converse inclusion. The correction explains how to separate an element outside the radical by a prime ideal, using a maximal ideal in a suitable nonzero localization of the quotient ring. The relevant background is \wikilink{s26--b006}{localization} and \wikilink{s28--b003}{contraction of prime ideals}. This general-ring argument is distinct from \wikilink{a2s26}{the Nullstellensatz over a field}.
\end{remark}
